\section{Experiments}\label{sec:experiments}
\subsection{Set-up, baselines, metrics and statistics}
Every controller runs in the same loop: the plant \eqref{eq:model} with the true parameters, the state delivered with the noise of Table~\ref{tab:noise} scaled by the noise level of the cell, the commands sampled every \SI{50}{ms} and rate-limited, \SI{20}{s} of simulation. Table~\ref{tab:controllers} lists the nine controllers. The threshold rule (THR) reproduces the autopilot logic; the fixed lift-sharing schedule (SCH) follows the nominal trim of the current airspeed with a pusher loop on the airspeed and a rotor loop on the flight-path rate; the gain-scheduled PID (GS-PID) schedules its gains on the airspeed; the model predictive controller (MPC) solves a nominal-model tracking problem on a \SI{4}{s} horizon by sequential conic programming; the feed-forward (MLP) and recurrent (LSTM) policies use the bounded head of HOT and are trained by the same physics-informed procedure; HOT$-$LSM is the transformer without the safety module; PPO~\cite{Schulman2017} and SAC~\cite{Haarnoja2018} are trained with Stable-Baselines3~\cite{Raffin2021} on the same plant family with the negative of \eqref{eq:loss} as the reward. Metrics: the time to the target airspeed (first entry within \SI{0.3}{m/s}), the energy of the identified power model, the minimum lift fraction, the maximum angle of attack, the extreme flight-path angles and the minimum sink rate, the percentage of samples with a violated constraint, the intervention rate of the safety module and the wall time per update. Statistics: mean and standard deviation over plants and seeds, the Wilcoxon signed-rank test at $p<0.05$ against HOT ($+/=/-$ counts of the cells in which HOT is better, equal or worse) and the Friedman mean rank~\cite{Demsar2006}.
\begin{table}[t]\centering\footnotesize
\caption{Controllers compared.}\label{tab:controllers}
\setlength{\tabcolsep}{4pt}
\begin{tabular}{llcc}\toprule
label & controller & learned & guarantee\\\midrule
THR & autopilot threshold rule & -- & --\\
SCH & fixed lift-sharing schedule & -- & --\\
GS-PID & gain-scheduled PID & -- & --\\
MPC & nominal-model MPC, \SI{4}{s} horizon & -- & nominal plant\\
MLP & feed-forward policy, bounded head & \checkmark & --\\
LSTM & recurrent policy on the window & \checkmark & --\\
HOT$-$LSM & transformer without the safety module & \checkmark & --\\
PPO, SAC & reinforcement learning, same reward & \checkmark & --\\
HOT & this article & \checkmark & Theorem~\ref{thm:main}\\
\bottomrule\end{tabular}\end{table}

\subsection{Training}\label{sec:trainres}
%% Fig: training curves (loss terms, violation rate, reached fraction) for HOT, MLP, LSTM; text with numbers
Fig.~\ref{fig:train} shows the training of HOT, the feed-forward and the recurrent policy under the same procedure: the time term, the barrier term, the fraction of episodes that reach the target within the horizon and the plant-head loss over the iterations. The barrier term falls by \trainBarrierDrop{} over the training and the reached fraction rises to \trainReached{} at the final horizon of \SI{15}{s}; the plant-head loss falls to \trainAux{}, i.e.\ the context token has learned to carry the plant. Training took \trainWallMin{}~min on four processor cores.
\begin{figure*}[t]\centering\includegraphics[width=\textwidth]{fig05_training.pdf}
\caption{Training curves under the physics-informed procedure: (a) time term; (b) barrier term; (c) fraction of episodes that reach the target within the curriculum horizon; (d) plant-head loss of HOT.}\label{fig:train}
\end{figure*}


\subsection{Nominal transitions}\label{sec:nominal}
%% Fig: nominal front and back transitions, all controllers (V, lift fraction, alpha, commands); Table: metrics
Fig.~\ref{fig:nominal} shows both transitions on the identified plant for every controller and Table~\ref{tab:nominal} the metrics. HOT reaches cruise in \hotFrontT{}~s and the hover-exit airspeed in \hotBackT{}~s with the lift fraction never below \hotFrontLf{} and \hotBackLf{}, the angle of attack below \hotFrontAl{}$^\circ$ and \hotBackAl{}$^\circ$, and \hotFrontEnergy{}~Wh and \hotBackEnergy{}~Wh of energy; the safety module intervened at \hotFrontInterv{} of the \hotFrontUpd{} front updates and \hotBackInterv{} of the \hotBackUpd{} back updates, and the same policy without the module, HOT$-$LSM, flies the same times (\hotnfFrontT{} and \hotnfBackT{}~s) but crosses the pitch box on the way. The threshold rule is fast in the front transition because it cuts the rotors as soon as the airspeed crosses its threshold, at the price of the lowest lift fraction of all controllers, \thrFrontLf{}; the schedule and the gain-scheduled PID keep the lift near unity but need \schFrontT{}~s and \pidFrontT{}~s in the front and 11 to 15~s in the back transition; the nominal-model MPC keeps the lift fraction at unity and takes \mpcFrontT{}~s with its \SI{4}{s} horizon and rate limits. The feed-forward and recurrent policies, run with the same module, take \mlpFrontT{}~s and \lstmFrontT{}~s in the front transition and hold less lift; SAC reaches the target in \sacFrontT{}~s and \sacBackT{}~s, and PPO does not reach cruise within the window on the identified plant, its rotor command oscillating between its bounds (Fig.~\ref{fig:nominal}d).
\begin{figure}[t]\centering\includegraphics[width=\columnwidth]{fig06_nominal.pdf}
\caption{Nominal front (left) and back (right) transitions: airspeed, lift fraction with the \SI{75}{\percent} bound, angle of attack with its bound, rotor thrust, for eight controllers (SAC, close to PPO, is in Table~\ref{tab:nominal}).}\label{fig:nominal}
\end{figure}
\begin{table*}[t]\centering\footnotesize
\caption{Metrics on the identified plant.}\label{tab:nominal}
\begin{tabular}{l|cccc|cccc}\toprule
& \multicolumn{4}{c|}{front transition ($6\to15$ m/s)} & \multicolumn{4}{c}{back transition ($15\to6$ m/s)}\\
controller & time [s] & energy [Wh] & min $\lf$ & max $\alpha$ [deg] & time [s] & energy [Wh] & min $\lf$ & max $\alpha$ [deg]\\\midrule
THR & 4.6 & 2.99 & 0.84 & -0.0 & 12.2 & 3.51 & 0.95 & 3.2\\
SCH & 7.2 & 3.38 & 0.99 & -0.0 & 11.3 & 3.77 & 0.98 & 0.8\\
GS-PID & 5.2 & 3.32 & 0.99 & -0.0 & 14.5 & 3.56 & 0.98 & -0.0\\
MPC & 18.4 & 3.39 & 1.00 & 0.4 & 19.5 & 3.60 & 1.00 & 0.3\\
MLP & 3.8 & 4.02 & 0.92 & -0.1 & 6.9 & 4.51 & 0.92 & 1.9\\
LSTM & 6.6 & 4.16 & 0.97 & -0.0 & 5.2 & 4.87 & 0.94 & -0.8\\
HOT$-$LSM & 4.0 & 4.17 & 0.94 & -0.2 & 4.7 & 4.51 & 0.96 & 0.6\\
PPO & -- & 8.94 & 0.32 & 1.2 & 8.2 & 4.03 & 0.94 & 3.2\\
SAC & 7.3 & 3.01 & 0.91 & -0.3 & 4.9 & 4.30 & 0.94 & 4.7\\
\textbf{HOT} & 4.5 & 4.21 & 0.94 & -0.2 & 4.7 & 4.51 & 0.96 & 0.6\\
\bottomrule\end{tabular}
\end{table*}


\subsection{Generalisation across the plant family}\label{sec:cross}
%% Fig: box plots per metric over the unseen plants, Friedman ranks; Table: Wilcoxon +/=/-
Table~\ref{tab:cross} and Fig.~\ref{fig:cross} report the \numCross{} unseen plants of the box, none of which was seen in training and none of which is known to any controller, with a front and a back transition each. HOT reaches the target band within the \SI{20}{s} window on \hotCrossFrontReached{} of the \hotCrossFrontN{} front and \hotCrossBackReached{} of the \hotCrossBackN{} back transitions, in a median of \hotCrossFrontTmed{}~s and \hotCrossBackTmed{}~s when it does, keeps the lift fraction above \hotCrossLf{} on every plant and violates no constraint on any of them; the module intervened at \hotCrossFrontIntervPct{}\% of the front and \hotCrossBackIntervPct{}\% of the back updates. The Wilcoxon test on the time to target is significant at $p<0.05$ in favour of HOT against every controller except its own unfiltered version, and the Friedman mean rank over the \numFriedmanN{} transitions flown by all controllers is \hotCrossRank{} for HOT against \bestOtherCrossRank{} for the best alternative, \bestOtherCrossName{}. The price of the guarantee is energy, about one watt-hour more than the classical controllers, spent by the rotors during the shorter transition, and the lift fraction, which the schedule and the gain-scheduled PID hold at 0.97 by taking twice the time. The transitions that miss the window fall into two groups. In the front transition the policy commands a nose-down attitude on strong-wing plants that cannot be certified against the light-wing plants still retained in the set, so the certified command holds the rotors and the aircraft climbs instead of accelerating; in the back transition \hotCrossBackStalledUnver{} plants stall after a single unverified backup step left the policy in a state it had not learned from, so it decelerated by drag alone and stopped just above the band, and \hotCrossBackStalledNoInt{} stall without any intervention. The transformer without the module, HOT$-$LSM, reaches \hotnfCrossReached{} of the transitions in \hotnfCrossT{}~s but violates a constraint on \hotnfCrossViol{}\% of the samples over the family, the threshold rule on \thrCrossViol{}\%, the schedule on \schCrossViol{}\% and PPO on \ppoCrossViol{}\%; the gain-scheduled PID and the MPC violate nothing but take \pidCrossT{} and \mpcCrossT{}~s, and the MPC reaches \mpcCrossReached{} of its transitions. The module is therefore what turns the fastest policy into a safe one, at the cost of the stalls just described.
\begin{table*}[t]\centering\footnotesize
\caption{Generalisation across the plant family: \numCross{} unseen plants, front and back transitions, mean $\pm$ standard deviation; W: Wilcoxon signed-rank test against HOT ($+$: HOT better, $=$: no difference, $-$: HOT worse, $p<0.05$); rank: Friedman mean rank on the time to target (lower is better).}\label{tab:cross}
\begin{tabular}{l|c|cc|cc|cc|c|c}\toprule
controller & reached & time [s] & W & energy [Wh] & W & min $\lf$ & W & viol.\ [\%] & rank\\\midrule
THR & 0.96 & 8.0$\pm$3.9 & $+$ & 3.18$\pm$0.67 & $-$ & 0.88$\pm$0.06 & $+$ & 1.6 & 3.82\\
SCH & 0.87 & 9.4$\pm$2.9 & $+$ & 3.62$\pm$0.44 & $-$ & 0.97$\pm$0.03 & $-$ & 2.5 & 5.24\\
GS-PID & 0.95 & 9.7$\pm$4.5 & $+$ & 3.50$\pm$0.34 & $-$ & 0.97$\pm$0.03 & $-$ & 0.0 & 5.06\\
MPC & 0.50 & 15.1$\pm$3.4 & $+$ & 3.44$\pm$0.22 & $-$ & 0.96$\pm$0.04 & $-$ & 0.0 & 7.28\\
HOT$-$LSM & 0.98 & 5.2$\pm$2.1 & $-$ & 4.39$\pm$0.28 & $=$ & 0.92$\pm$0.07 & $=$ & 1.6 & 1.89\\
PPO & 0.50 & 8.3$\pm$1.1 & $+$ & 6.48$\pm$2.49 & $+$ & 0.60$\pm$0.31 & $+$ & 36.6 & 5.61\\
SAC & 0.87 & 5.9$\pm$2.0 & $+$ & 3.66$\pm$0.67 & $-$ & 0.91$\pm$0.03 & $+$ & 0.3 & 3.86\\
\textbf{HOT} & 0.89 & 5.5$\pm$2.5 &  & 4.38$\pm$0.25 &  & 0.94$\pm$0.02 &  & 0.0 & 3.24\\
\bottomrule\end{tabular}
\end{table*}
\begin{figure*}[t]\centering\includegraphics[width=\textwidth]{fig07_cross.pdf}
\caption{Distribution over the \numCross{} unseen plants of the time to target, the energy, the minimum lift fraction and the maximum angle of attack for each controller (boxes: quartiles; whiskers: extremes; dotted: bounds).}\label{fig:cross}
\end{figure*}


\subsection{Outside the box, gusts and noise}\label{sec:robust}
Table~\ref{tab:robust} and Fig.~\ref{fig:sweeps} extend the test beyond the training distribution: \numOutbox{} plants drawn from a box enlarged by \SI{50}{\percent} about its centre, vertical gusts of 0.25 to \SI{1.5}{m/s} standard deviation and sensor noise of half to twice the identified level, ten seeds each. Outside the box HOT reaches the band on \hotOutFrontReached{} of \hotOutFrontN{} front and \hotOutBackReached{} of \hotOutBackN{} back transitions with the lift fraction above \hotOutLf{} and no violation; the module intervenes at \hotOutFrontIntervPct{}\% of the front updates, and \hotOutFrontStalledNoInt{} front and \hotOutBackStalledNoInt{} back misses happen without any intervention, which is the policy itself leaving its training distribution. The threshold rule reaches \thrOutReached{} of these transitions with the lift fraction at \thrOutLf{} and \thrOutViol{}\% of the samples in violation. Gusts are the weak point of the certificate: the backup is rolled out without the gust, and a gust moves the flight path by more than the measurement bound $\varepsilon$ in one step, so the consistency test discards plants for the wrong reason and the module falls back on the unverified backup on \hotGustStrongUnverShare{}\% of its interventions at 1 and \SI{1.5}{m/s} (\hotGustStrongInterv{}\% of the updates). At \SI{0.25}{m/s} nothing is violated; at 1 and \SI{1.5}{m/s} the angle-of-attack and lift bounds are crossed on \hotGustStrongViol{}\% of the samples, against \pidGustOneViol{} and \pidGustOneHalfViol{}\% for the gain-scheduled PID, \thrGustOneHalfViol{}\% for the threshold rule and \sacGustOneHalfViol{}\% for SAC at \SI{1.5}{m/s}, while HOT keeps the highest minimum lift fraction of all controllers under the strongest gust, \hotGustOneHalfLf{} against \thrGustOneHalfLf{} and \sacGustOneHalfLf{}. Enlarging $\varepsilon$ by the gust bound restores the coverage of Lemma~\ref{lem:cons} at the price of a larger retained set. Under sensor noise the module keeps every constraint at every level; the back transition is unaffected (\hotNoiseTwoBackReached{} of 10 within the window at twice the noise, median \hotNoiseTwoBackTmed{}~s), whereas the front transition slows with the noise, median \hotNoiseOneFrontTmed{}~s at the identified level and \hotNoiseTwoFrontTmed{}~s with \hotNoiseTwoFrontReached{} of 10 within the window at twice the level: the policy's commands lose their coherence in the noisy window, with the module intervening at only \hotNoiseInterv{}\% of the updates.
\begin{table*}[t]\centering\footnotesize
\caption{Outside the box, gusts and noise: fraction of transitions that reach the target within \SI{20}{s}, minimum lift fraction (mean $\pm$ standard deviation) and fraction of samples in violation of a constraint.}\label{tab:robust}
\begin{tabular}{l|ccc|ccc|ccc}\toprule
& \multicolumn{3}{c|}{outside the box} & \multicolumn{3}{c|}{gusts} & \multicolumn{3}{c}{noise}\\
controller & reached & min $\lf$ & viol.\ [\%] & reached & min $\lf$ & viol.\ [\%] & reached & min $\lf$ & viol.\ [\%]\\\midrule
THR & 0.89 & 0.87$\pm$0.07 & 1.3 & 0.96 & 0.78$\pm$0.17 & 9.3 & 1.00 & 0.89$\pm$0.06 & 0.0\\
SCH & 0.78 & 0.95$\pm$0.05 & 4.3 & 0.96 & 0.83$\pm$0.15 & 9.8 & 1.00 & 0.97$\pm$0.01 & 0.0\\
GS-PID & 0.78 & 0.95$\pm$0.05 & 0.0 & 0.84 & 0.84$\pm$0.14 & 3.6 & 1.00 & 0.97$\pm$0.01 & 0.0\\
PPO & 0.48 & -3.31$\pm$19.95 & 38.4 & 0.50 & 0.54$\pm$0.30 & 37.4 & 0.50 & 0.62$\pm$0.30 & 35.9\\
SAC & 0.84 & 0.90$\pm$0.04 & 0.3 & 1.00 & 0.79$\pm$0.16 & 6.0 & 1.00 & 0.91$\pm$0.02 & 0.1\\
\textbf{HOT} & 0.66 & 0.92$\pm$0.05 & 0.0 & 0.89 & 0.87$\pm$0.08 & 8.4 & 0.87 & 0.95$\pm$0.01 & 0.0\\
\bottomrule\end{tabular}
\end{table*}
\begin{figure*}[t]\centering\includegraphics[width=\textwidth]{fig09_sweeps.pdf}
\caption{Time to target and minimum lift fraction against the gust standard deviation (a, b) and the sensor-noise level (c, d), mean and standard deviation over ten seeds and both transitions.}\label{fig:sweeps}
\end{figure*}


\subsection{From the states recorded in flight}\label{sec:replay}
The dataset~\cite{Dony2026data} records the aircraft's own hand-offs under the threshold rule, including a front transition in which the rotors were cut early and the wing lost lift for three seconds, and a back transition that ran to a high angle of attack. Fig.~\ref{fig:replay} starts HOT from the state recorded \SI{2.7}{s} before that rotor cut, at \SI{5.4}{m/s} and climbing, and from the state recorded as the back transition began, at \SI{13.4}{m/s} with the rotors off. Both states lie outside the training distribution of trims. From the front state HOT keeps the lift fraction above \hotReplayFrontLf{} where the recorded flight lost its lift, but the module holds the rotors through \hotReplayFrontInterv{} certified interventions in the low-speed climb, where the plant set is least contracted, and cruise is not reached within the \SI{20}{s} window; from the back state it spins the rotors up within the step bound, keeps the angle of attack below \hotReplayBackAl{}$^\circ$, where the recorded flight went to $36^\circ$, and reaches the band in \hotReplayBackT{}~s with \hotReplayBackInterv{} interventions.
\begin{figure}[t]\centering\includegraphics[width=\columnwidth]{fig10_replay.pdf}
\caption{Angle of attack and lift fraction of HOT started from the states recorded in flight (solid) beside the recorded continuation under the threshold rule (dashed), front (left) and back (right).}\label{fig:replay}
\end{figure}


\subsection{PX4 firmware in the loop}\label{sec:sil}
Fig.~\ref{fig:sil} shows both transitions with the PX4 firmware in the loop, from the trims reached under the same commands. The front transition reaches cruise in \silFrontT{}~s and the back transition enters the hover-exit band at \silBackT{}~s, both slower than on the identified model (Section~\ref{sec:nominal}): the firmware's aircraft carries the drag of its fuselage and tail plates and a slower attitude loop than the identified one, and the policy, which is not told which plant it is flying, commands the pitch and the pusher it learned for the box, so the airspeed follows the drag of the aircraft it meets; in the back transition the last metre per second is bled by drag alone, with the pusher at zero and the rotors near hover thrust, which takes ten seconds. The angle of attack stays within $[\silFrontAlMin^\circ,\silFrontAlMax^\circ]$ and $[\silBackAlMin^\circ,\silBackAlMax^\circ]$, the flight path within $\silFrontGam^\circ$ and $\silBackGam^\circ$, and the lift fraction evaluated on the identified model never falls below \silFrontLf{} and \silBackLf{}. The pitch setpoint carries the noise of the firmware's airspeed estimate, which the attitude loop filters out of the pitch; the module certified every command (\silFrontInterv{} interventions in \silFrontUpd{} and \silBackInterv{} in \silBackUpd{} updates). The update, network and check together, took \silUpdMs{}~ms on average in the loop, on a processor shared with the parallel evaluations.
\begin{figure*}[t]\centering\includegraphics[width=0.72\textwidth]{fig11_sil.pdf}
\caption{Closed loop with the PX4 firmware in software-in-the-loop, front (left) and back (right) transitions: airspeed from the autopilot's estimator, pitch setpoint sent to the attitude loop with the pitch and angle of attack it produced, rotor and pusher commands.}\label{fig:sil}
\end{figure*}


\subsection{Ablations and parameter analysis}\label{sec:ablation}
Fig.~\ref{fig:ablation} evaluates the architecture on the first \numAblPlants{} plants of the cross-plant cell, both transitions, with the bare policy and with the module: HOT with one module removed, the feed-forward and recurrent policies trained by the same procedure, and reduced sizes. The reduced variants and the module ablations were trained for 500 iterations, HOT, MLP and LSTM for 800. Without the physics features the bare policy reaches the band on \ablNoPhysReached{}\% of the transitions against \ablFullReached{}\% for HOT and violates on \ablNoPhysViol{}\% of the samples; without the auxiliary head the context token still resolves the parameters (mean $R^2$ \ablNoAuxRtwo{} against \interpRtwoMean{}), which the pooled physics features carry on their own, and the policy reaches \ablNoAuxReached{}\%. The feed-forward policy trained by the same physics-informed procedure, with the same physics features and previous command, also flies the box (\ablMlpReached{}\% of the transitions in a median of \ablMlpTmed{}~s), and the recurrent policy \ablLstmReached{}\% in \ablLstmTmed{}~s: the training procedure and the physics features of Section~\ref{sec:method} transfer to every architecture, and the history-conditioned transformer adds to them the plant estimate of Section~\ref{sec:interp}, which the feed-forward policy cannot form. With the module every variant is free of violations. A window of 8 samples, a width of 32 and a single block reach \ablLsmallReached{}, \ablDsmallReached{} and \ablBoneReached{}\% of the transitions bare at the shorter training budget.
\begin{figure*}[t]\centering\includegraphics[width=0.8\textwidth]{fig12_ablation.pdf}
\caption{Ablations and parameter analysis on \numAblPlants{} unseen plants, both transitions: (a) fraction of transitions that reach the band within \SI{20}{s}, bare policy and with the module, for HOT, HOT without the physics features, without the auxiliary head, the LSTM and the MLP; (b) fraction of samples in violation for the bare policies; (c) reach fraction of the reduced sizes (window 8, width 32, one block; trained for 500 iterations).}\label{fig:ablation}
\end{figure*}


\subsection{Interpretability: attention and the context token}\label{sec:interp}
Fig.~\ref{fig:interp} looks inside the plant-inference attention. The attention of the context token over the window is nearly uniform, between \interpAttMin{} and \interpAttMax{} against $1/L=\interpAttUniform{}$ (a): the token pools the whole window rather than a few samples, and the plant information sits in the pooled physics features. Panels (b) and (c) follow the $R^2$ between the context-token estimate and the true parameter along the transition on the unseen plants: each parameter is resolved while the response that reveals it is inside the window and released once the window has slid past it. The lift-curve offset is read off the wing-borne trim at the start of the back transition ($R^2=\interpRtwoCLzero{}$ at \interpTCLzero{}~s), the rotor gain once the rotors respond in the \interpTrBetaR{} transition ($\interpRtwoBetaR{}$ at \interpTBetaR{}~s), the attitude-loop rate from the first pitch response in the \interpTrKappa{} transition ($\interpRtwoKappa{}$ at \interpTKappa{}~s) and the lift-curve slope during the acceleration ($\interpRtwoCLa{}$ at \interpTCLa{}~s); the drag parameters, which act on the airspeed slowly, are not resolved ($R^2$ at most \interpRtwoCDzero{}). The estimates are shrunk towards the centre of the box by the measurement noise, to \interpShrinkCLzero{} of the true spread for $C_{L0}$, as a least-squares head under noise should be, so (d) and (e) show them standardised. A sliding window implies this forgetting; the ablation of the auxiliary head (Section~\ref{sec:ablation}) measures what the estimate is worth to the policy.
\begin{figure*}[t]\centering\includegraphics[width=\textwidth]{fig13_interp.pdf}
\caption{Interpretability on the unseen plants: (a) attention of the context token over the window, front and back, against the uniform weight $1/L$; (b), (c) $R^2$ of the context-token estimate of the lift-curve offset, the lift-curve slope, the attitude-loop rate and the rotor gain against the true values along the front and the back transition; (d), (e) standardised estimates against the standardised true values at the time of the best $R^2$.}\label{fig:interp}
\end{figure*}


\subsection{Time cost}\label{sec:time}
Table~\ref{tab:time} lists the update times on one core of the development computer in interpreted Python. The network evaluates in \hotEvalMs{}~ms; the safety check takes \hotCheckMs{}~ms on average and \hotCheckMaxMs{}~ms at most with the online plant set, most of it in the rollouts of the backup, and it is skipped when the network command has already been certified at the previous step for a state that has not moved; the nominal-model MPC needs \mpcStepMs{}~ms per update.
\begin{table}[t]\centering\footnotesize
\caption{Properties of transition controllers and learned controllers with safety (\checkmark: yes, --: no).}\label{tab:compare}
\setlength{\tabcolsep}{2.5pt}
\begin{tabular}{lcccccc}\toprule
& \rotatebox{90}{learned policy} & \rotatebox{90}{plant family} & \rotatebox{90}{implicit identification} & \rotatebox{90}{lift guarantee} & \rotatebox{90}{lagged actuators} & \rotatebox{90}{flight-identified}\\\midrule
Threshold rule~\cite{PX4docs} & -- & -- & -- & -- & -- & \checkmark\\
Gain-scheduled / hierarchical~\cite{Lyu2017,Verling2016} & -- & -- & -- & -- & -- & \checkmark\\
Optimised manoeuvre~\cite{Maqsood2012} & -- & -- & -- & -- & -- & --\\
Transition MPC~\cite{Li2018,Allenspach2020,Bauersfeld2021} & -- & -- & -- & -- & -- & \checkmark\\
Adaptive VTOL control~\cite{Dong2020} & -- & \checkmark & -- & -- & -- & --\\
Approximate MPC~\cite{Hertneck2018,Karg2020,Chen2018} & \checkmark & -- & -- & -- & -- & --\\
Predictive / backup safety filters~\cite{Wabersich2021,Gurriet2020} & \checkmark & -- & -- & -- & -- & --\\
Domain randomisation, adaptation~\cite{Peng2018,Kumar2021} & \checkmark & \checkmark & \checkmark & -- & -- & --\\
This article & \checkmark & \checkmark & \checkmark & \checkmark & \checkmark & \checkmark\\
\bottomrule\end{tabular}\end{table}
\begin{table}[t]\centering\footnotesize
\caption{Update time per controller (one core, interpreted code).}\label{tab:time}
\begin{tabular}{lc}\toprule
controller & time per update\\\midrule
threshold rule & \thrUs{}~\si{\micro\second}\\
schedule, GS-PID & \schUs{}~\si{\micro\second}, \pidUs{}~\si{\micro\second}\\
MPC (sequential conic programming) & \mpcStepMs{}~ms\\
HOT network & \hotEvalMs{}~ms\\
HOT safety check (mean, max) & \hotCheckMs{}, \hotCheckMaxMs{}~ms\\
\bottomrule\end{tabular}\end{table}


\subsection{Comparison with previous work}\label{sec:compare}
Table~\ref{tab:compare} places HOT against the transition controllers and the learned controllers of Section~\ref{sec:related} on the properties that this article set out to combine. The autopilot rule and the gain-scheduled and optimised transitions of the VTOL literature plan nothing online and carry no guarantee on the lift; model predictive transition controllers plan online but on one nominal model; approximate model predictive control and the safety-filter literature certify a learned or a nominal command on one plant; domain-randomised policies generalise across plants but without a constraint guarantee. HOT is the only entry that learns a single policy over a plant family, infers the plant it is flying from its own history, guarantees the lift margin, the angle of attack, the flight path and the sink rate for every plant of the verified set with lagged actuators and the attitude loop inside the plant, and is validated on a flight-identified aircraft with the autopilot firmware in the loop.
\begin{table}[t]\centering\footnotesize
\caption{Properties of transition controllers and learned controllers with safety (\checkmark: yes, --: no).}\label{tab:compare}
\setlength{\tabcolsep}{2.5pt}
\begin{tabular}{lcccccc}\toprule
& \rotatebox{90}{learned policy} & \rotatebox{90}{plant family} & \rotatebox{90}{implicit identification} & \rotatebox{90}{lift guarantee} & \rotatebox{90}{lagged actuators} & \rotatebox{90}{flight-identified}\\\midrule
Threshold rule~\cite{PX4docs} & -- & -- & -- & -- & -- & \checkmark\\
Gain-scheduled / hierarchical~\cite{Lyu2017,Verling2016} & -- & -- & -- & -- & -- & \checkmark\\
Optimised manoeuvre~\cite{Maqsood2012} & -- & -- & -- & -- & -- & --\\
Transition MPC~\cite{Li2018,Allenspach2020,Bauersfeld2021} & -- & -- & -- & -- & -- & \checkmark\\
Adaptive VTOL control~\cite{Dong2020} & -- & \checkmark & -- & -- & -- & --\\
Approximate MPC~\cite{Hertneck2018,Karg2020,Chen2018} & \checkmark & -- & -- & -- & -- & --\\
Predictive / backup safety filters~\cite{Wabersich2021,Gurriet2020} & \checkmark & -- & -- & -- & -- & --\\
Domain randomisation, adaptation~\cite{Peng2018,Kumar2021} & \checkmark & \checkmark & \checkmark & -- & -- & --\\
This article & \checkmark & \checkmark & \checkmark & \checkmark & \checkmark & \checkmark\\
\bottomrule\end{tabular}\end{table}

